% MCFG #2 v72: window--finite-abelian-group extension, using untyped context shortening.
\section{Windows combined with finite abelian groups}
\label{app:window-abelian}


\paragraph{Semigroup convention (v82 migration).}
The group morphism $g:\Sigma^*\to A$ is retained as a computational
monoid representation. The learning bias is the \emph{positive}
semigroup restriction of the window--group product, with a fresh
identity adjoined only for evaluating empty template factors.
By Lemma~\ref{lem:semigroup-transfer}, nonempty component tests and
the selected characteristic words are unchanged. The finite-group
zero-sum argument uses $g$ only as an auxiliary invariant.


This appendix extends the polynomial total-length theorem to
finite-abelian-group refinements of a prefix--suffix window.
The original Theorem~\ref{thm:boundary-polydata},
which does not involve a group, remains unchanged.

\begin{lemma}[Short tuples with prescribed abelian-group values]
\label{lem:group-typed-shortest}
Let $G=(V,\Sigma,P,S,\fan)$ be a reduced good
MCFG of fan-out at most $f$ and rule rank at most
$r_{\max}\ge1$. Let
$g:\Sigma^*\to A$ be a homomorphism into a fixed
finite abelian group. Write $\mathsf d(H)$ for the
maximal length of a zero-sum-free sequence over a
finite abelian group $H$, so that $\mathsf D(H)=\mathsf d(H)+1$
is its Davenport constant.
For every $B\in V$ and
every realizable vector
$\mathbf s\in A^{\fan(B)}$, there is
$\tuple u\in L_B(G)$ with
$g^{(\fan(B))}(\tuple u)=\mathbf s$
and
\(\|\tuple u\|_1\le\beta_g:= (r_{\max}+1)\bar\theta_G \max\{1,\mathsf d(A^f)\,|V||A|^f\}.\)
\end{lemma}

\begin{proof}
Choose a terminal derivation $T$ from $B$ whose
root has $g$-vector $\mathbf s$, minimizing
first the total yield length and then the number
of derivation nodes. Label each node by its
nonterminal $C$ and the componentwise $g$-vector
of its subtree yield. There are at most
$n_g:=|V||A|^f$ possible such labels.

If an ancestor and a proper descendant had
the same label, replacing the former subtree
by the latter would preserve the root
$g$-vector: every MCFG rule computes the
parent's group values from the values of its
children. The total yield length cannot
increase because the rules are nondeleting,
and the number of nodes decreases.
This contradicts the choice of $T$.
Thus its depth is at most $n_g$.

Call a node \emph{heavy} if the total length
of its subtree yield exceeds $\bar\theta_G$.
All ancestors of a heavy node are heavy
by nondeletion. Suppose $v_1,\ldots,v_m$
are pairwise incomparable heavy nodes and
the root has fan-out $d\le f$.
For each $v_i$ labelled by $C_i$,
choose a terminal derivation from $C_i$
of total length at most $\bar\theta_G$,
which exists by
Lemma~\ref{lem:thickness-local-bounds}(ii).
Replacing the subtree at $v_i$ by this
short derivation strictly decreases the
total yield length.

Because the grammar is linear and nondeleting,
the output components of mutually incomparable
subtrees contribute disjoint terminal blocks
to the root components, possibly in different
orders. In an abelian group, multiplication
is additive and the order of those blocks is
immaterial. Therefore there are vectors
$\delta_i\in A^d$ such that simultaneous
replacement of any subset $J$ of the $v_i$
changes the root $g$-vector by
$\sum_{i\in J}\delta_i$.
If $m>\mathsf d(A^d)$, the sequence
$\delta_1\cdots\delta_m$ over $A^d$ has a nonempty
zero-sum subsequence $(\delta_i)_{i\in J}$, so
simultaneous replacement at the nodes $v_i$,
$i\in J$, preserves the root
vector while strictly shortening its yield.
This contradicts minimality. Every heavy
antichain consequently has at most
$\mathsf d(A^d)\le\mathsf d(A^f)$ nodes; the
inequality holds because $A^d$ embeds in $A^f$.

If the root is light, its length is at most
$\bar\theta_G$. Otherwise the heavy nodes
form an ancestor-closed tree. Its leaves
are a heavy antichain, and every heavy node
belongs to a path to one of these leaves.
There are therefore at most
$\mathsf d(A^f)\,n_g$ heavy nodes.
Every explicit terminal factor written by
a rule is of total length at most
$\bar\theta_G$ by
Lemma~\ref{lem:thickness-local-bounds}(i).
Every light child of a heavy node contributes
at most $\bar\theta_G$ additional terminals,
and a heavy node has at most $r_{\max}$
children. Counting the heavy nodes and their
light children gives the asserted bound.
The case $|A|=1$ is covered by the light-root
alternative.
\end{proof}

\begin{remark}[Size of the zero-sum constant]
\label{rem:davenport-size}
The prefix-sum pigeonhole argument gives
$\mathsf d(A^f)\le|A|^f-1$, so
Lemma~\ref{lem:group-typed-shortest} is never weaker
than the bound obtained from consecutive partial sums.
It is usually much stronger.  Write
$A\cong C_{n_1}\oplus\cdots\oplus C_{n_s}$ and
$n=\exp(A)\ge2$. Then
\[
 f\sum_{i=1}^{s}(n_i-1)\ \le\ \mathsf d(A^f)\ \le\
 (n-1)+n\log\frac{|A|^f}{n},
\]
where the lower bound is attained by a basis
sequence and the upper bound is due to van Emde Boas
and Kruyswijk~\cite[Proposition~5.1.8 and Theorem~5.5.5]{GeroldingerHalterKoch2006}.
If $A$ is a $p$-group, the lower bound is an
equality~\cite[Theorem~5.5.9]{GeroldingerHalterKoch2006}
(Olson's theorem).
Hence $\mathsf d(A^f)=O(\exp(A)\,f\log|A|)$ grows
linearly in the fan-out bound, whereas $|A|^f-1$ grows
exponentially; for $A=C_2$ the antichain bound drops
from $2^f-1$ to~$f$.

The value $\mathsf d(A^d)$ is also the limit of the
antichain step as stated.  That step uses only the
fact that the difference vectors $\delta_i$ have no
nonempty zero-sum subsequence, and a zero-sum-free
sequence of length $\mathsf d(A^d)$ exists.  We do not
claim that a grammar realizing such an extremal
configuration exists.
\end{remark}

\begin{lemma}[Short canonical tuples for window--abelian typing]
\label{lem:window-abelian-short}
Assume the hypotheses of
Lemma~\ref{lem:group-typed-shortest} and fix
$k,\ell\ge0$ and a finite typing
$s:\Sigma^+\to\mathsf S$ satisfying
$\widehat{s}\preceq(\widehat{s}_{k,\ell},g)$.
Let $\widetilde G_0$ be the trimmed
full $\widehat{s}$-refinement of $G$, and let
$n_{\mathrm{typ}}$ be its number of
non-start nonterminals. For each such
nonterminal $X$ one may choose a
canonical tuple $\omega(X)$ with
\(\|\omega(X)\|_1\le f(k+\ell+1)n_{\mathrm{typ}}\, \bigl(\bar\theta_G+r_{\max}\beta_g\bigr).\)
\end{lemma}

\begin{proof}
Choose, among the typed terminal derivations
from $X$, one with the fewest nodes.
A repeated typed nonterminal on a path
could be eliminated while preserving the
root $\widehat{s}$-type, so each path has at most
$n_{\mathrm{typ}}$ nodes.
As in the proof of
Lemma~\ref{lem:boundary-short-canonical},
mark in each root component the positions
that determine its $\widehat{s}_{k,\ell}$-type:
all positions if the component is shorter
than the window threshold, otherwise the
first $k$ and last $\ell$ positions, with
one position marked when $k=\ell=0$.
The ancestor paths to all marks form a
skeleton with at most
$f(k+\ell+1)n_{\mathrm{typ}}$ nodes.

Each child subtree omitted from this
skeleton lies wholly inside the unmarked
interior of every root component to which
it contributes. Record its nonterminal
$C$ and its actual componentwise $g$-value
$\mathbf s$. Replace it by a derivation
from $C$ with the same $\mathbf s$
and length at most $\beta_g$, as supplied
by Lemma~\ref{lem:group-typed-shortest}.
Goodness keeps all replacement components
nonempty. No protected boundary position
changes, and each long root component
remains long, so its window type is
unchanged. Moreover, each root component
is a concatenation of the explicit
terminal material and the child components;
each replaced child component retains its
$g$-value, so the root $g$-vector is also
unchanged, without requiring a
commutativity argument at this stage.

The product $(\widehat{s}_{k,\ell},g)$ and hence
the coarser $\widehat{s}$-type of the root are
unchanged. Forget the intermediate typed
labels and uniquely lift the new
underlying $G$-derivation back to the
$\widehat{s}$-refinement using
Proposition~\ref{prop:output-invariants}(ii).
Because $X$ is retained, the resulting
root derivation survives trimming.
Each skeleton node writes at most
$\bar\theta_G$ terminals and has at most
$r_{\max}$ omitted children. The
displayed length estimate follows.
\end{proof}

\begin{theorem}[Polynomial data for window--abelian typing]
\label{thm:window-abelian-polydata}
Fix $f,r_{\max},k,\ell,\Sigma$, a
finite abelian group $A$, a homomorphism
$g:\Sigma^*\to A$, and a finite typing
$\widehat{s}\preceq(\widehat{s}_{k,\ell},g)$.
For every $L\in\Cmcf{f}{s}$ and reduced good
fan-out-$\le f$, rank-$\le r_{\max}$ grammar $G$
with $L(G)=L\cap\Sigma^+$, the capped
reconstruction operator $\mathcal B^{\le r_{\max}}_{f,s}$
has a characteristic sample
satisfying
\[
 |\Wit(\widetilde G_0)|=O(\|G\|),\qquad
 \|\Wit(\widetilde G_0)\|_+
   =O_{f,r_{\max},k,\ell,s,g,\Sigma}
     (\|G\|^3\bar\theta_G).
\]
\end{theorem}

\begin{proof}
The fixed finite typing gives
$n_{\mathrm{typ}}\le|V||M|^f=O(\|G\|)$.
Lemma~\ref{lem:group-typed-shortest}
gives $\beta_g=O(\|G\|\bar\theta_G)$.
By Lemma~\ref{lem:window-abelian-short},
the selected canonical tuple lengths are
$O(\|G\|^2\bar\theta_G)$.
Lemma~\ref{lem:untyped-contexts}
therefore bounds each characteristic
witness word by $O(\|G\|^2\bar\theta_G)$;
Proposition~\ref{prop:witness-cardinality}
bounds their number by $O(\|G\|)$.
The $\eps$ datum, if any, adds only a
constant to the encoded total.
\end{proof}

\begin{remark}
Commutativity is used in the zero-sum
antichain step of
Lemma~\ref{lem:group-typed-shortest}.
At fan-out one, subtrees contribute
single intervals, and the companion CFG
paper treats arbitrary finite groups.
For $f\ge2$ and nonabelian groups,
the pieces from different subtrees can
interleave in different root components.
The zero-sum argument in this appendix does not cover
that case.  The separate order-homogenization proof in
Appendix~\ref{app:window-group} does cover arbitrary
finite groups, including nonabelian groups.
\end{remark}
